\subsection{Reducing the fit cost when few coefficients are free}
\label{sec:slice_models}

Most results are likelihood scans in one or two Wilson coefficients, with the rest held
at their SM values or a few of them profiled. The workspace, however, carries all forty
templates of the basis and evaluates all forty terms of Eq.~\eqref{eq:smeftpara_reorg}
in every bin, including the terms that depend only on coefficients the fit never moves.
Those terms are constants, and once they are constant the weight $w_i(c)$ of many basis
templates is exactly zero everywhere in the scan. A template whose weight is zero
contributes nothing to any prediction the fit will make, and can be left out of the
workspace.

One free coefficient leaves three of the forty terms and two leave six, and the number of
templates with a non-vanishing weight follows: between five and fifteen on a single axis,
about ten on a plane, and all forty once every coefficient moves. The count depends on
the direction and not only on the number of free coefficients --- the
$c_{\PH\mathrm{kin}}$ axis keeps six templates and the $c_{\PH G}$ axis fifteen --- and
it follows from the basis in one evaluation for any other scan. The coefficients that
have to be kept are those the fit scans, profiles, or holds at a value other than the SM
one.

Nothing new has to be simulated. The smaller model is the same card with templates
removed, so it can be built for any scan from the basis already in hand.

What is saved is the size of the statistical model. The expected yield of a bin is
$\sum_i w_i(c) N_{ij}$, so the work per bin is the number of templates times the number
of terms, and the two shrink together; Table~\ref{tab:slices} gives both, for the full
basis and for two scans, with the cost of one prediction and of one scan point. The
scan-point column is measured on a complete analysis model --- of order twenty
categories, parametric shapes with discrete profiling, the full set of systematic
uncertainties, and the $80$ templates of the nine-coefficient basis at two
centre-of-mass energies --- in which one point of a nine-coefficient scan takes about
half an hour against the tens of seconds an analysis with a handful of signal templates
is used to. Those are order-of-magnitude figures from one model and will differ between
analyses.

A global fit floats all nine coefficients, so it carries every template and saves
nothing. Three precautions apply to it, as general practice rather than as responses to
an observed failure. The minimiser is run with tight settings, since what a loose
tolerance saves is small next to a scan point that has not converged. Several starting
points are tried, because the cross section is quadratic in the coefficients and the
likelihood need not be unimodal: in a realistic nine-coefficient scan the $c_{\PQt G}$
direction shows a second minimum near $c_{\PQt G}=1$. And a coefficient the selection
does not respond to is fixed rather than floated, with the choice stated, so that no flat
direction enters the fit. On the measurement of
Sec.~\ref{sec:validation_subleading_role} those are $c_{QQ\PQt\PQt}$ and
$c_{QQ}^{(8)}$, whose acceptance varies by less than $0.1$ over their full range in
either $4\PQb$ selection: they move the rate and leave the shape untouched, so a fit that
floats them is fitting a direction its selection cannot see.
\textcolor{red}{Whether that survives at a collider reaching much higher \mHH, where the
four-top operators grow fastest, is a question for the $100$~TeV study
(Sec.~\ref{sec:energy_beyond}).}

For a channel as light as the \HHbbbb example used here the difference changes little in
practice: the model has few categories, and a scan that takes half an hour per point is
still affordable. The reduction is worth making where the model is heavy --- many
categories, several decay channels combined into one fit, and scans that have to be
repeated for each variation --- which is where the cost of carrying forty templates
through every bin becomes the limiting one.

The numerical results are unchanged, as they have to be: the templates that are dropped
carry a weight of zero, and the agreements measured on the \HHbbbb example are in
Table~\ref{tab:slices}, down to $10^{-5}$ in $-2\Delta\ln L$ with up to three
coefficients profiled and $2\times10^{-3}$ with up to seven. The Monte Carlo statistical
uncertainty of the morphed prediction is unchanged for a separate reason: the templates
are reweightings of the same simulated events, and their morphing reproduces the direct
reweighting to the target point in any basis (Sec.~\ref{sec:uncertainties}). Measured
along the $c_{\PH\mathrm{kin}}$ axis with the full covariance of the templates, the two
models give the same $\sigma_{\text{MM}}$ to the three digits quoted, from $0.173\%$ of
the yield at the SM point to $0.240\%$ at $c_{\PH\mathrm{kin}}=8$; integrated over \mHH
they agree to $0.1\%$ and per bin to within $4\%$. The interpolation uncertainty belongs
to the input samples rather than to the basis.

\begin{table}[htbp]
\centering
\begin{tabular}{lcccccc}
\hline
free & terms & templates & products & direct & one scan & agreement with \\
coefficients & kept & kept & per bin & evaluation & point & the full model \\
\hline
all nine (full basis)             & $40$ & $40$ & $1600$ & $9\,\mu$s   & $\sim30$~min & reference \\
$c_{\PH}$, $c_{\PH\mathrm{kin}}$  & $6$  & $10$ & $60$   & ---         & ---          & $5\times10^{-7}$ \\
$c_{\PH\mathrm{kin}}$             & $3$  & $6$  & $18$   & $0.2\,\mu$s & $\sim6$~min  & $<10^{-9}$ \\
\hline
\end{tabular}
\caption{Cost and agreement of the reduced model. ``Products per bin'' is the number of
templates times the number of terms kept. The direct evaluation is the prediction of the
\HHbbbb example outside any fitting framework; the scan point is one point of a
likelihood scan in a complete analysis model and is an order of magnitude only.
Agreement is the largest difference in $-2\Delta\ln L$ between the reduced and the full
model over the scans tested.}
\label{tab:slices}
\end{table}

A smaller model again can be built by giving the scan its own basis points --- three on
one axis, six on a plane, as many as the terms it keeps --- instead of keeping whichever
of the forty templates survive there. Those points are reweighted templates the card
does not already contain, so this is worthwhile only when the same scan is repeated.

Nothing in the construction is specific to the SMEFT parameterization: it needs only
that the prediction is a polynomial and that the morphing weights are polynomials in the
same terms. The same arguments apply to any physics case built with the matrix method,
HEFT among them.
